\section{Related Systems Survey}
\subsection{Canvas LMS}
\setcounter{table}{0}
\renewcommand{\thetable}{2.1.\arabic{table}}
\begin{xltabular}{\textwidth}{>{\bfseries}p{2.8cm} X X}
    % 1. Header for the first page
    \caption{Comparative Analysis of Academic LMS Platforms} \label{tab:lms_comparison} \\
    \toprule
    Feature / Criterion & Moodle & Canvas LMS \\
    \midrule
    \endfirsthead

    % 2. Repeating header for all subsequent pages
    \multicolumn{3}{l}{\textit{\small Table \thetable\ (continued)}} \\
    \toprule
    Feature / Criterion & Moodle & Canvas LMS \\
    \midrule
    \endhead

    % 3. Footer on all pages except the last
    \midrule
    \multicolumn{3}{r}{\textit{\small (Continued on next page)}} \\
    \endfoot

    % 4. Footer for the final page
    \bottomrule
    \endlastfoot

    % ----------------- TABLE CONTENT -----------------
    Target Environment & Higher education, K-12, and enterprise training. & Primarily higher education and K-12. \\
    \addlinespace
    System Architecture & Open-source. Highly modular, allowing deep database access and custom plugin development. & Proprietary (SaaS). Cloud-hosted with strict core codebase protection, relying on external LTI integrations. \\
    \addlinespace
    Course Registration & Managed via standard enrollment plugins or manual cohort syncs. Lacks native AI-driven prerequisite validation and advisement. & Integrates with university Student Information Systems (SIS). Does not feature an intelligent AI advisor for course selection. \\
    \addlinespace
    Course \& Document Management & Extensive file support. Instructors can upload syllabi and materials. & Excellent module-based organization and drag-and-drop file management. \\
    \addlinespace
    Assessment \& Grading & Comprehensive quizzing engine with rubrics and automated grading. & Highly intuitive ``SpeedGrader'' interface with built-in rubrics and annotation tools. \\
    \addlinespace
    Early Warning \& Analytics & Rule-based activity completion tracking and standard analytics models to flag inactive students. & Strong native learning analytics displaying student engagement, grade distribution, and interaction metrics. \\
    \addlinespace
    AI Integration & No native AI out-of-the-box. However, its open architecture makes it the perfect foundation for building custom modules (like your document-grounded Q\&A or risk-detection engines). & Relies on third-party integrations (e.g., external AI tutors via LTI). Cannot easily modify the core system to embed a native AI registration advisor. \\
    \addlinespace
    Role \& User Management & Extremely granular Role-Based Access Control (RBAC). Admins can define custom permissions down to specific activity levels. & Standardized roles (Admin, Teacher, TA, Student) with solid, but less granular, permission customization. \\
\end{xltabular}

\noindent Traditional LMS platforms like Moodle and Canvas function primarily as passive content repositories. The proposed system transitions into an active, AI-driven educational environment by natively embedding intelligence into university workflows. It surpasses standard systems through two core capabilities:

\begin{itemize}
    \item \textbf{Contextual AI Assistant (RAG):} Embeds a Q\&A chatbot directly grounded in uploaded course materials, eliminating reliance on fragmented third-party plugins.
    \item \textbf{Diagnostic Risk Detection:} Proactively identifies at-risk students and provides explicit causal explanations (e.g., consecutive scoring downtrends) rather than just displaying raw statistical charts.
\end{itemize}


\subsection{D2L Brightspace}

\begin{xltabular}{\textwidth}{>{\bfseries}p{2.8cm} X X}

    % 1. Header for the first page
    \caption{Comparative Analysis of D2L Brightspace}
    \label{tab:brightspace_comparison} \\

    \toprule
    Feature / Criterion & Moodle & D2L Brightspace \\
    \midrule
    \endfirsthead

    % 2. Repeating header for all subsequent pages
    \multicolumn{3}{l}{\textit{\small Table \thetable\ (continued)}} \\
    \toprule
    Feature / Criterion & Moodle & D2L Brightspace \\
    \midrule
    \endhead

    % 3. Footer on all pages except the last
    \midrule
    \multicolumn{3}{r}{\textit{\small (Continued on next page)}} \\
    \endfoot

    % 4. Footer for the final page
    \bottomrule
    \endlastfoot

    % ----------------- TABLE CONTENT -----------------

    Target Environment &
    Supports higher education, K--12 education, and organizational
    training through a flexible learning management platform. &
    Supports higher education, K--12 education, and organizational
    learning through an integrated learning management platform. \\

    \addlinespace

    System Architecture &
    Open-source and modular, allowing institutions to customize the
    platform and develop or install plugins. &
    Commercial LMS platform with integrated learning tools and
    additional analytics capabilities available through its product
    offerings. \\

    \addlinespace

    Course Registration &
    Supports course enrollment through configurable enrollment methods,
    plugins, and administrative settings. &
    Supports course management and enrollment workflows. Its Student
    Success System focuses on predicting learner success rather than
    providing AI-based course recommendations. \\

    \addlinespace

    Course \& Document Management &
    Supports course creation, learning resources, activities, and
    file uploads through its course management features. &
    Supports course content organization, learning activities, and
    resources within the Brightspace learning environment. \\

    \addlinespace

    Assessment \& Grading &
    Provides quizzes, assignments, grading workflows, and configurable
    assessment activities. &
    Provides assessment and grading functions integrated with the
    learning environment. Its Student Success System also uses
    assessment-related indicators for predictive analysis. \\

    \addlinespace

    Early Warning \& Analytics &
    Provides activity completion tracking, reports, and learning
    analytics capabilities. Additional functionality can depend on
    installed plugins and institutional configuration. &
    The Student Success System uses predictive analytics and machine
    learning techniques to estimate learner success and predicted
    grades, helping instructors identify students who may need support. \\

    \addlinespace

    AI Integration & 
    Moodle's extensible architecture supports integration and
    customization through plugins and external services. Available
    AI capabilities depend on the implementation and installed tools. &
    The Student Success System integrates predictive analytics and
    machine learning into learner success monitoring, using historical
    course data to generate predictions and support early intervention. \\


    \addlinespace

    Role \& User Management &
    Supports role-based permissions for administrators, teachers, and
    students. Institutions can customize access and extend functionality
    through plugins and configuration. &
    Provides role-based access for administrators, instructors, and
    learners, with permissions managed through the Brightspace
    administrative environment. \\



\end{xltabular}

\noindent Moodle and D2L Brightspace both provide core learning management capabilities, including course management, learning resources, and assessment support. Their approaches to customization and learning analytics differ. Moodle's open-source and modular architecture allows institutions to adapt the platform through plugins and integrations, while D2L Brightspace's Student Success System provides predictive analytics and machine learning capabilities to estimate learner success using historical course data. 

\noindent This comparison highlights two relevant considerations for the project: the flexibility to integrate and customize LMS functionality, as demonstrated by Moodle, and the use of predictive analytics to support student monitoring, as demonstrated by D2L Brightspace. These considerations can inform the selection of the system architecture and the integration of AI capabilities according to the project's requirements.




